\HLASourcePage{17}{16}
\HLAAnchor{clause:1}
\pdfbookmark[0]{1. Overview}{bookmark:clause:1}
\thispagestyle{hla-opening}
\begingroup
\setlength{\parskip}{0pt}
\vspace*{-15.5pt}
{\sffamily\bfseries\fontsize{24pt}{26pt}\selectfont
  IEEE Standard for\\
  Modeling and Simulation (M\&S)\\
  High Level Architecture (HLA)---\\
  Federate Interface Specification\par
}
\vspace{27.1pt}
\HLAClause{1. Overview}
\vspace{18.5pt}
\HLAAnchor{clause:1.1}\pdfbookmark[1]{1.1 Scope}{bookmark:clause:1.1}\HLASubclause{1.1 Scope}
\vspace{0.8pt}
\endgroup

This document defines the interface between federates (simulations, supporting utilities, or interfaces to live
systems) and the underlying software services that support interfederate communication in a distributed
simulation domain.

\begingroup\setlength{\parskip}{0pt}
\vspace{0.4pt}
\HLAAnchor{clause:1.2}\pdfbookmark[1]{1.2 Purpose}{bookmark:clause:1.2}\HLASubclause{1.2 Purpose}
\vspace{9.7pt}
\endgroup

The High Level Architecture (HLA) has been developed to provide a common architecture for distributed
modeling and simulation (M\&S). To facilitate interfederate communications, HLA federates interact with an
underlying software infrastructure. This specification defines the standard services and interfaces to be used
by the federates to support efficient information exchange when participating in a distributed federation
execution. (A federation execution occurs when sets of federates are brought together to support an
objective.)

\begingroup\setlength{\parskip}{0pt}
\vspace{0.4pt}
\HLAAnchor{clause:1.3}\pdfbookmark[1]{1.3 Word usage}{bookmark:clause:1.3}\HLASubclause{1.3 Word usage}
\vspace{9.7pt}
\endgroup

The word \emph{shall} indicates mandatory requirements strictly to be followed in order to conform to the standard
and from which no deviation is permitted (shall equals is required to).\textsuperscript{\fontsize{7.98pt}{8pt}\selectfont 6,7}

The word \emph{should} indicates that among several possibilities one is recommended as particularly suitable,
without mentioning or excluding others; or that a certain course of action is preferred but not necessarily
required (should equals is recommended that).

The word \emph{may} is used to indicate a course of action permissible within the limits of the standard (may equals
is permitted to).

The word \emph{can} is used for statements of possibility and capability, whether material, physical, or causal (can
equals is able to).

\begingroup\setlength{\parskip}{0pt}
\vspace{0.4pt}
\HLAAnchor{clause:1.4}\pdfbookmark[1]{1.4 Background}{bookmark:clause:1.4}\HLASubclause{1.4 Background}
\vspace{9.7pt}
\endgroup

This document is being developed to define the HLA interface specification. The HLA requires a language-
independent specification (LIS) and multiple language bindings to support intersimulation communication
in a distributed simulation domain.

\vspace{0.3pt}
\HLACompactFootnote{6}{The use of the word \emph{must} is deprecated and cannot be used when stating mandatory requirements, \emph{must} is used only to describe
unavoidable situations.}
\HLACompactFootnote{7}{The use of the word \emph{will} is deprecated and cannot be used when stating mandatory requirements, \emph{will} is only used in statements of fact.}
\clearpage
